\documentclass[border=6pt]{standalone}
\usepackage{fontspec}
\usepackage{xeCJK}
\usepackage{xcolor}
\usepackage{tikz}
\usetikzlibrary{
    arrows.meta,
    backgrounds,
    calc,
    decorations.pathreplacing,
    fit,
    positioning,
    shapes.geometric
}

\setmainfont{Times New Roman}
\setsansfont{Noto Sans SC}
\setCJKmainfont{Noto Sans SC}
\setCJKsansfont{Noto Sans SC}

\definecolor{ink}{HTML}{1F2937}
\definecolor{muted}{HTML}{64748B}
\definecolor{lineblue}{HTML}{2563EB}
\definecolor{prefill}{HTML}{E7F0FF}
\definecolor{prefillLine}{HTML}{2F5597}
\definecolor{decode}{HTML}{EAF7EA}
\definecolor{decodeLine}{HTML}{2F7D32}
\definecolor{kv}{HTML}{FFF3D8}
\definecolor{kvLine}{HTML}{B7791F}
\definecolor{accent}{HTML}{FDE8E8}
\definecolor{accentLine}{HTML}{C2410C}
\definecolor{panel}{HTML}{F8FAFC}
\definecolor{panelLine}{HTML}{94A3B8}
\definecolor{violet}{HTML}{F1ECFF}
\definecolor{violetLine}{HTML}{6D28D9}
\definecolor{teal}{HTML}{DFF7F3}
\definecolor{tealLine}{HTML}{0F766E}

\tikzset{
    every picture/.style={font=\sffamily},
    title/.style={font=\sffamily\bfseries\Large, text=ink},
    subtitle/.style={font=\sffamily\small, text=muted, align=center},
    block/.style={
        draw=panelLine,
        rounded corners=5pt,
        fill=white,
        line width=0.8pt,
        align=center,
        inner xsep=7pt,
        inner ysep=6pt
    },
    pblock/.style={block, fill=prefill, draw=prefillLine},
    dblock/.style={block, fill=decode, draw=decodeLine},
    kvblock/.style={block, fill=kv, draw=kvLine},
    warnblock/.style={block, fill=accent, draw=accentLine},
    vblock/.style={block, fill=violet, draw=violetLine},
    tblock/.style={block, fill=teal, draw=tealLine},
    cluster/.style={
        draw=panelLine,
        rounded corners=6pt,
        fill=panel,
        line width=0.8pt,
        inner sep=8pt
    },
    arrow/.style={-{Latex[length=3mm,width=2mm]}, draw=ink, line width=0.85pt},
    bluearrow/.style={arrow, draw=prefillLine},
    greenarrow/.style={arrow, draw=decodeLine},
    kvarrow/.style={arrow, draw=kvLine},
    orangearrow/.style={arrow, draw=accentLine},
    control/.style={-{Latex[length=2.6mm,width=1.8mm]}, draw=muted, line width=0.75pt, dashed},
    label/.style={font=\sffamily\scriptsize, text=muted, align=center},
    small/.style={font=\sffamily\scriptsize, align=center},
    rank/.style={
        draw=panelLine,
        rounded corners=3pt,
        fill=white,
        minimum width=12mm,
        minimum height=6mm,
        font=\sffamily\scriptsize,
        align=center
    },
    head/.style={
        draw=ink,
        line width=0.35pt,
        minimum width=6.5mm,
        minimum height=4.8mm,
        font=\sffamily\tiny,
        align=center
    },
    metric/.style={
        draw=panelLine,
        rounded corners=4pt,
        fill=white,
        font=\sffamily\scriptsize,
        align=center,
        inner xsep=5pt,
        inner ysep=4pt
    }
}


\begin{document}
\begin{tikzpicture}[x=1cm,y=1cm]
    \node[title] at (8.5,8.15) {PD 分离中的 KV Cache 布局映射与转换};
    \node[subtitle] at (8.5,7.72) {P/D 两侧 TP、PP 与注意力结构可能不同，Connector 需要把源端 KV 精确转换成 Decode 侧 PageAttention 可消费的 block 布局};

    \node[pblock, minimum width=4.7cm, minimum height=5.6cm] (prefill) at (2.85,4.25) {};
    \node[font=\sffamily\bfseries\small, text=prefillLine] at (2.85,6.65) {Prefill 侧 KV 生产};
    \node[small, text=muted] at (2.85,6.25) {示例：P TP=4，PP=2};

    \foreach \y/\stage in {5.45/PP stage 0,4.6/PP stage 0,3.35/PP stage 1,2.5/PP stage 1} {
        \node[rank, minimum width=1.2cm] at (1.08,\y) {\stage};
    }
    \foreach \i/\y in {0/5.45,1/4.6,2/3.35,3/2.5} {
        \node[rank, fill=white, minimum width=1.0cm] at (2.15,\y) {P\i};
        \node[head, fill=teal] at (3.1,\y) {$h_{\i}$};
        \node[head, fill=kv] at (3.78,\y) {$B_{0..n}$};
    }
    \node[small, text=muted] at (2.85,1.75) {GQA/MHA/MQA：KV head 随 TP rank 切分};

    \node[kvblock, minimum width=4.25cm, minimum height=5.6cm] (connector) at (8.45,4.25) {};
    \node[font=\sffamily\bfseries\small, text=kvLine] at (8.45,6.65) {KV Connector};
    \node[metric, fill=white, minimum width=3.25cm] (s1) at (8.45,5.72) {1. 感知 TP/PP/layer partition};
    \node[metric, fill=white, minimum width=3.25cm] (s2) at (8.45,4.9) {2. 选择 P rank 或 rank 集合};
    \node[metric, fill=white, minimum width=3.25cm] (s3) at (8.45,4.08) {3. HCCL/HCCS/RDMA 传输};
    \node[metric, fill=white, minimum width=3.25cm] (s4) at (8.45,3.26) {4. split / cat / transpose by block};
    \node[metric, fill=white, minimum width=3.25cm] (s5) at (8.45,2.44) {\_npu\_page\_load + reshape\_and\_cache};
    \node[small, text=accentLine] at (8.45,1.55) {目标：减少逐层逐 block 的 copy/reshape 与 kernel launch};

    \node[dblock, minimum width=4.75cm, minimum height=5.6cm] (decode) at (14.0,4.25) {};
    \node[font=\sffamily\bfseries\small, text=decodeLine] at (14.0,6.65) {Decode 侧 KV 消费};
    \node[small, text=muted] at (14.0,6.25) {PageAttention block table / slot mapping};
    \foreach \r/\x in {0/12.75,1/15.25} {
        \node[rank, fill=white, minimum width=1.45cm] at (\x,5.42) {D rank \r};
    }
    \foreach \x/\h in {12.25/h_0,12.82/h_1,14.72/h_2,15.29/h_3} {
        \node[head, fill=teal] at (\x,4.82) {$\h$};
    }
    \foreach \i/\x in {0/12.18,1/12.75,2/13.32,3/14.62,4/15.19,5/15.76} {
        \node[head, fill=kv, minimum width=5.3mm] at (\x,4.15) {$B_{\i}$};
    }
    \node[metric, fill=decode, minimum width=3.4cm] at (14.0,3.25) {本地 KV block 写入};
    \node[metric, fill=white, minimum width=3.4cm] at (14.0,2.45) {Decode forward 读取分片 KV};
    \node[small, text=decodeLine] at (14.0,1.75) {保证输出一致性，并在 D 写入完成后释放 P 侧缓存};

    \draw[kvarrow] (prefill.east) -- node[label, above] {源端 KV layout} (connector.west);
    \draw[kvarrow] (connector.east) -- node[label, above] {目标端 block layout} (decode.west);
    \draw[control] (s1.east) to[out=0,in=180] (decode.170);
    \draw[control] (s2.west) to[out=180,in=0] (prefill.10);

    \node[warnblock, minimum width=4.55cm, minimum height=0.75cm] at (3.55,0.6) {MLA：可在等价 P rank 中负载均衡};
    \node[warnblock, minimum width=8.25cm, minimum height=0.75cm] at (10.9,0.6) {GQA/MHA/MQA：必须按 KV head 和 layer 精确匹配，避免 head 缺失或 block 无法消费};
\end{tikzpicture}
\end{document}
